% Solutions/hints for ch16 exercises (assembled into Appendix C)

\begin{solution}{exr:ch16-resolution}
(a)~Take the wall $[4.00,4.04]\times[0,10]$ and the horizontal segment from $\vect{a}=(3.995,5)$ to $\vect{b}=(4.045,5)$. Its length is $0.05=\delta$, so the check uses only the two endpoints; both lie outside the wall ($3.995<4.00$ and $4.045>4.04$) and the segment crosses it. With the inflation by $\delta/2=0.025$ the point $\vect{a}$, at distance $0.005$ from the wall, is rejected. (b)~Let $\vect{o}$ be an obstacle point and write the squared distance from $\vect{o}$ to the point at arc length $s$ along a segment as $g(s)=(s-s^{*})^{2}+h_{*}^{2}$, where $h_{*}$ is the distance from $\vect{o}$ to the line and $s^{*}$ the arc length of its projection. If some point $\vect{q}$ of the segment, between the test points at $s=0$ and $s=\delta$, has $\dist(\vect{q},\vect{o})=h\le r$, then $h_{*}\le h$ and one of the two test points has $(s-s^{*})^{2}\le\delta^{2}/4$ (the one on the far side of $s^{*}$, or either one if $s^{*}$ is outside $[0,\delta]$), so its distance to $\vect{o}$ is at most $\sqrt{h^{2}+\delta^{2}/4}\le\sqrt{r^{2}+\delta^{2}/4}$ and it fails a test with that margin. Hence the margin $\sqrt{r^{2}+\delta^{2}/4}$, which is smaller than $r+\delta/2$ for $r>0$, also guarantees a collision-free segment. It is tight: a point obstacle at perpendicular distance exactly $r$ from the midpoint between two test points is at distance $\sqrt{r^{2}+\delta^{2}/4}$ from both of them, so any smaller margin accepts the segment although the robot touches the obstacle at the midpoint. (c)~The ball of radius $d_0$ around $\vect{a}$ contains no point within $r$ of an obstacle, so every segment point closer than $d_0$ to $\vect{a}$ is free without a test. The adaptive check therefore jumps to the point at distance $d_0$ along the segment, computes its own free distance $d_1$, jumps by $d_1$, and so on, until it passes $\vect{b}$ (the segment is free) or until a free distance is $\le0$ (collision). In open space a few jumps cover the whole segment; near an obstacle the jumps shrink, and an implementation stops and reports a collision once a jump falls below a minimum length, which keeps the check conservative and finite. This is the ``conservative advancement'' scheme used by many collision libraries.
\end{solution}

\begin{solution}{exr:ch16-completeness}
(a)~Call iteration $i$ a success if its sample lands in the ball $B_{K_{i-1}+1}$, the one just beyond the stage reached so far. The ball is determined by the past and the sample is independent of the past and uniform with probability $1-p_{\mathrm{goal}}$, so the conditional success probability is at least $p$ whatever happened before; by a standard coupling the number of successes $S_n$ in $n$ iterations is stochastically at least $\mathrm{Bin}(n,p)$. Every success raises the stage by at least one (the claim in the proof of \cref{thm:ch16-complete}), so $T>n$ implies $S_n<m$ and $\Prob(T>n)\le\Prob(\mathrm{Bin}(n,p)<m)$. For $n>2m/p$ the mean $np$ exceeds $2m$, and the Chernoff bound $\Prob(\mathrm{Bin}(n,p)\le(1-\alpha)np)\le\exp(-\alpha^{2}np/2)$ with $\alpha=1/2$ gives $\Prob(T>n)\le\exp(-np/8)$, an exponential decay in $n$. (b)~Take \cref{ex:ch16-map} with $p_{\mathrm{goal}}=1$. Every sample is the goal, so every iteration extends the vertex nearest to $(9,9)$ by a step of $\eta$ along the straight line towards it; after four steps the segment enters the wall $[3,5]\times[0,6]$, is rejected, and the same rejected extension is repeated for ever, although a path exists. In the proof, the probability that a sample lands in a ball off the straight line is $p=(1-p_{\mathrm{goal}})\,\mathrm{vol}(B)/\mathrm{vol}(\mathcal{X})=0$. (c)~The proof uses $\eta>0$ once, to choose $\nu\le\eta/3$ so that a sample within $3\nu$ of the tree is reached exactly by \textsc{Steer}. With $\eta_i=1/i$ the balls must shrink with $i$, their volumes are proportional to $i^{-d}$, and $\sum_i i^{-d}<\infty$ for $d\ge2$: the expected total number of samples that ever land in the current target ball is finite, so with positive probability the tree stops advancing after finitely many stages and the algorithm is not probabilistically complete. (The tree may still reach the goal by other routes, but the argument no longer proves that it does.)
\end{solution}

\begin{solution}{exr:ch16-voronoi}
The bisector of $\vect{v}_0$ and $\vect{v}_1$ is the line $x=0.5$ and that of $\vect{v}_1$ and $\vect{v}_2$ the line $x=1.5$, so the Voronoi regions inside $\mathcal{X}=[-1,3]\times[-2,2]$ (area $16$) are the strips $-1\le x\le0.5$ (area $6$), $0.5\le x\le1.5$ (area $4$) and $1.5\le x\le3$ (area $6$). The next sample selects $\vect{v}_0$ with probability $6/16=0.375$, $\vect{v}_1$ with $0.25$ and $\vect{v}_2$ with $0.375$: the interior vertex is the least likely, and the root is a tip of the chain too. For a chain of $k+1$ vertices spaced $\eta$ apart and centred in a square of side $L\gg k\eta$, the bisectors of consecutive vertices are parallel lines $\eta$ apart, so every interior vertex owns a strip of width $\eta$ and area at most $\eta L$, probability at most $\eta/L$; the two tips share the rest, $(L^{2}-(k-1)\eta L)/2$ each, so each tip is chosen with probability $\tfrac12-(k-1)\eta/(2L)$, close to $\tfrac12$. The chain therefore grows almost only at its two ends, and mostly at the end that faces the larger empty region. Centring matters: if the chain starts near a corner and points into the map, as in \cref{ex:ch16-map}, the region of the single tip is the part of the square beyond the last bisector, which is almost all of it, so that tip is selected with probability close to $1$ rather than $1/2$.
\end{solution}

\begin{solution}{exr:ch16-goalbias}
Expect a U-shaped curve: with $p_{\mathrm{goal}}=0$ the tree spends its samples on the whole map, while with a large bias almost every sample is the goal, so the vertex nearest the goal is extended again and again into the wall and the tree stops growing elsewhere (\cref{sec:ch16-pitfalls}); the median is smallest in between, near $p_{\mathrm{goal}}=0.05$--$0.1$, matching \cref{tab:ch16-connect}, where $0.5$ is worse than no bias at all. Moving the goal to $(9,1)$ removes one of the two walls from the straight line between start and goal, so greedy steps towards the goal are blocked over a shorter stretch and recover sooner: the penalty for a large bias shrinks and the best bias moves upwards.
\end{solution}

\begin{solution}{exr:ch16-shortcut}
(a)~The replaced part of the path is a polyline from $\vect{a}$ to $\vect{b}$, and the triangle inequality (applied along it) gives its length as at least $\norm{\vect{a}-\vect{b}}$, the length of the new segment; the rest of the path is unchanged. The new path is collision-free because the untouched segments were, and the new segment was accepted by \textsc{CollisionFree}. (b)~\textsc{Prune} outputs only input vertices, in strictly increasing index order, starting at $\state_0$ and ending at $\state_k$ (the jump to $j=i+1$ is always free), so the output has at most $k+1$ vertices. Example: $\state_0=(0,0)$, $\state_1=(1,1.2)$, $\state_2=(2,0)$, $\state_3=(3,3)$ with a disc obstacle of radius $0.05$ at $(0.5,0.5)$. The disc blocks $[\state_0,\state_3]$ (the diagonal passes through its centre) and nothing else, since it is at distance $0.5$ from $[\state_0,\state_2]$, about $0.064$ from $[\state_0,\state_1]$ and further from the other segments. Greedy pruning from $\state_0$ tries $\state_3$ (blocked) and jumps to $\state_2$, then to $\state_3$: length $2+\sqrt{10}\approx5.16$. The path $\state_0,\state_1,\state_3$ is also free and has length $\sqrt{2.44}+\sqrt{7.24}\approx4.25$. (c)~Every point of the initial path $(1,5),(1,9.75),(9,9.75),(9,5)$ has $y\ge5$. A shortcut replaces a sub-path by a segment between two of its points, and every point of that segment is a convex combination of two points with $y\ge5$, so the invariant $y\ge5$ survives every shortcut. Passing below the wall $[4,6]\times[2,9.5]$ requires a point with $4\le x\le6$ and $y<2$, which never appears; the paths obtained by shortcutting all cross the wall's column above it and are at least as long as $(1,5),(4,9.5),(6,9.5),(9,5)$, that is $2\sqrt{9+20.25}+2\approx12.82$, whereas the route below the wall through $(4,2)$ and $(6,2)$ has length $2\sqrt{9+9}+2\approx10.49$.
\end{solution}

\begin{solution}{exr:ch16-gridsizes}
The volume is $200\times200\times60$~m, so at $0.25$~m the lattice has
$800\times800\times240=1.536\times10^{8}$ cells, $153.6$~MB at one byte per
cell, and every interior cell has $3^{3}-1=26$ neighbours in the
26-connected lattice. Adding the velocity, $\pm4$~m/s per axis in steps of
$0.25$~m/s, gives $33$ values per axis and $33^{3}=35\,937$ velocity cells
per position cell: $1.536\times10^{8}\times35\,937\approx5.5\times10^{12}$
states, $5.5$~TB at one byte each, and $3^{6}-1=728$ neighbours if all axes
may change at once (in practice only the accelerations that respect the
dynamics are edges, which is fewer but still a lattice one cannot store).
An \rrt with $20\,000$ vertices in three dimensions needs
$20\,000\times(3\times8+4)=560\,000$~bytes, about $0.56$~MB: seven orders of
magnitude less than the velocity-augmented grid, because it stores only the
states it actually visited.
\end{solution}

\begin{solution}{exr:ch16-trace}
Iteration~11, $\state_{\mathrm{rand}}=(9.70,5.16)$: the distances to vertices~0
to~6 are smallest for vertex~6 at $(2.86,2.92)$, $\norm{\state_{\mathrm{rand}}-\state_6}=7.197$
(vertex~5 gives $7.424$), so $\state_{\mathrm{near}}=\state_6$. The unit vector
towards the sample is $(6.84,2.24)/7.197=(0.9504,0.3112)$, so
$\state_{\mathrm{new}}=\state_6+0.5\,(0.9504,0.3112)=(3.335,3.076)$, which is inside
the wall $[3,5]\times[0,6]$: the extension is rejected. Iteration~12,
$\state_{\mathrm{rand}}=(6.23,7.77)$: again vertex~6 is nearest ($5.906$ against
$6.355$ for vertex~5), the unit vector is $(0.5706,0.8212)$ and
$\state_{\mathrm{new}}=(3.145,3.331)$, again inside the wall and rejected. These
two lines are computed from the two-decimal coordinates of
\cref{tab:ch16-trace}; the trace printed by \code{ch16\_rrt.py}, which carries
the full precision of $\state_6$, gives $(3.332,3.073)$ and $(3.143,3.327)$ -- the
same conclusion, the difference being the rounding of the tabulated vertex. The Voronoi picture
explains the pattern: the tree is a chain that stops at $x\approx2.9$, so
the bisectors between consecutive vertices are almost vertical lines and the
region of the tip, vertex~6, contains the whole upper right of the map,
while vertex~5 owns the lower right. Every sample with $x>3$, about $70\,\%$
of the map, therefore selects vertex~5 or~6, and every step from them
towards such a sample crosses the wall.
\end{solution}
